\section{引言}\label{intro}

关于如何量子化一个经典辛流形，并不存在一个完全令人满意的一般理论，因为从经典力学到量子力学的过渡存在\textbf{反常}。量子算符的换位子与相空间上经典函数的 Poisson 括号并不一致，除非只取 $\hbar$ 的最低阶 \cite{G46,V51,Todorov}。

实践中，量子化通常是这样做的：粗略地说，把相空间坐标分成位置与动量，再定义位置函数的 Hilbert 空间。在几何量子化 \cite{Ko,So,Wo,Snia,Hall} 中，这一过程被形式化为相空间的一个“\textbf{极化}”的选取。所得量子理论一般确实依赖于极化的选取。在不同极化确实给出等价量子理论的特殊情形中，这往往会导致重要的结果。

几何量子化的一个替代方案是\textbf{膜量子化}。它是在 \cite{GW} 中提出的，其先导发展包括：二维拓扑场论的 $A$-模型与形变量子化之间关系的种种迹象 \cite{KO,BrS,Ka,Pe,Gu}，以及对若干有启发的例子的分析 \cite{AZ}。形变量子化 \cite{Many,WL,Fedosov,Kon,CF,CF2,SL} 是 $A$-模型故事的重要组成部分，它是量子化的一个近亲：给一个经典相空间配上一个量子形变的函数代数，但并不给出该代数在其上作用的 Hilbert 空间。按照 \cite{Fedosov} 的表述，形变量子化涉及的是辛联络的选取，而不是极化的选取。

在膜量子化中，起点不是经典相空间 $M$ 的极化的选取，而是 $M$ 的一个合适的\textbf{复化} $Y$ 的选取。$Y$ 应当是一个复辛流形，带有一个反全纯对合，使 $M$ 成为其不动点集的一个分支；此外 $Y$ 还应当有良好定义的 $A$-模型。于是 $M$ 的量子 Hilbert 空间可以定义为 $(\B,\B_\cc)$ 弦的 $A$-模型态空间，其中 $\B$ 是支撑在 $M$ 上的 Lagrangian $A$-膜，$\B_\cc$ 是 $A$-模型的某个\textbf{余各向同性膜}。余各向同性膜最初由 \cite{KO} 引入，并在后来的若干文章中得到了研究 \cite{Ka,Pe,Gu,AZ,KW,Herbst}。

在基于极化 $\P$ 的 $M$ 的几何量子化中，能够自然量子化为量子算符的 $M$ 上的函数，是那些通过 Poisson 括号生成保持极化的 $M$ 的典范变换的函数。在利用 $M$ 的复化 $Y$ 的膜量子化中，能够自然量子化的函数，是那些可以解析延拓为 $Y$ 上全纯函数的函数。通常 $Y$ 不是紧的，因此还要对函数在 $Y$ 的无穷远处的增长附加一个条件。

我们之所以重新考虑膜量子化，是因为涉及几何 Langlands 对应解析版本的最新进展 \cite{EFK,EFK2,EFK3,T}。正如我们在别处所述 \cite{GaW}，膜量子化是把几何 Langlands 的规范场论进路 \cite{KW} 推广到这一主题的解析版本时所需要的额外要素。

本文第 \ref{two} 节简要回顾建立量子化一般理论所遇到的困难，然后简要回顾几何量子化与形变量子化，再更详细地描述膜量子化的想法。有些内容我们讲得比 \cite{GW} 更精确，但在某些主题上，读者会在那篇文章中找到更多细节。

第 \ref{comparing} 节探讨几何量子化与膜量子化之间的关系。如果同一个相空间 $M$ 既可以用极化 $\P$ 作几何量子化，又可以用 $M$ 的复化 $Y$ 作膜量子化，那么能否对 $\P$ 与 $Y$ 给出条件，使得两种量子化方法一致？我们给出使这一点成立的条件；粗略地说，$\P$ 应当解析延拓为 $Y$ 的一个全纯极化 $\varPi$。

第 \ref{symq} 节探讨用 $M$ 与自身的 Lagrangian 对应、而不是 $M$ 上的函数，来定义膜量子化所得 Hilbert 空间上的算符。Lagrangian 对应就是 $M\times M$ 的一个 Lagrangian 子流形（两个因子上的辛结构相同而反号）。用对应而不是函数来定义算符，至少有两个理由。第一，这为研究膜量子化的对称性提供了一个有效的框架。第二，在一些重要的例子中，包括与几何 Langlands 相关的情形，$M$ 的复化 $Y$ 只有相对很少的全纯函数，取而代之的是与对应相关联的非常重要的算符。

第 \ref{qcomp} 节考虑一个特殊情形：把复辛流形 $Y$ 视为实辛流形来量子化。这一情形在几何 Langlands 的一个应用 \cite{GaW} 中很重要。

\begin{conceptbox}{膜量子化的三个物件}
膜量子化一旦选定复化 $Y$，剩下的就三样东西：
\begin{itemize}[leftmargin=1.2em,itemsep=0.15em,topsep=2pt]
\item 一个支撑在 $M$ 上的 Lagrangian $A$-膜 $\B$（$M$ 是反全纯对合 $\tau$ 的不动点集的一个分支）；
\item 一个支撑在全部 $Y$ 上的余各向同性膜 $\B_\cc$（“典范余各向同性膜”，其 Chan--Paton 丛在本文约定下平凡）；
\item 二者之间的弦态空间 $\Hom(\B,\B_\cc)$，它就是 $M$ 的量子 Hilbert 空间。
\end{itemize}
\end{conceptbox}

\begin{yizhu}[“反常”在这里指什么]
此处所说的反常，就是 Groenewold--Van Hove 意义下的\textbf{算符排序障碍}：不存在从经典多项式函数代数到算符代数的线性映射，使 Poisson 括号在 $\hbar$ 的所有阶都对应换位子。原文引用的 \cite{G46,V51} 正是这一 obstruction 的原始文献；\cite{Todorov} 是综述。[来源：A 级]
\end{yizhu}
